% Opening results and conventions, to follow V4's Motivation slide.
% Numerical provenance and qualifications: SOURCES.md and review/.
% Generated by questions_audit/make_slide4_no_cut.py: no GNN score cut; reco >=2 tracks.
\newcommand{\gnnEfficiency}{$68.68\pm0.06\%$}
\newcommand{\gnnPurity}{$69.45\pm0.06\%$}
\newcommand{\gnnAssigned}{$98.89\pm0.01\%$}
\newcommand{\gnnCleanReco}{$51.78\pm0.13\%$}
\newcommand{\amvfEfficiency}{$68.32\pm0.06\%$}
\newcommand{\amvfPurity}{$68.65\pm0.06\%$}
\newcommand{\amvfAssigned}{$99.52\pm0.01\%$}
\newcommand{\amvfCleanReco}{$35.17\pm0.10\%$}
\newcommand{\gnnGainEfficiency}{$+0.37\pm0.02$}
\newcommand{\gnnGainPurity}{$+0.80\pm0.02$}

\newcommand{\sourcecaption}[1]{\par\vspace{.5em}{\scriptsize\color{black!55}#1\par}}
\newcommand{\itkpaper}{\href{https://arxiv.org/html/2412.15090v2\#S5.SS5}{ITk performance paper, Sec.\ 5.5 (local draft pp.\ 27--29)}}
\newcommand{\runthreepaper}{\href{https://arxiv.org/abs/2308.09471}{Run 3 tracking paper, Fig.\ 13}}
\newcommand{\amvfnote}{\href{https://cds.cern.ch/record/2670380}{AMVF note (2019), Sec.\ 3.2, pp.\ 3--4}}

\begin{frame}{PV-Finder and AMVF: vertex finding}
\context{$14$ TeV $t\bar t$ simulation, $\langle\mu\rangle=192.5$, $|\eta|<4$; same 1,920 events; integral cut $0.3$.}
\centering
\small\renewcommand{\arraystretch}{1.2}
\begin{tabular}{lccc}
\toprule
Measurement & $z$ window & PV-Finder & AMVF\\
\midrule
All-vertex positional efficiency & $\pm0.1$ mm & $74.11\pm0.10\%$ & $70.99\pm0.10\%$\\
Unmatched reco / event & $\pm0.1$ mm & $25.16\pm0.11$ & $25.46\pm0.11$\\
Extra nearby reco / event & $\pm0.1$ mm & $0.112\pm0.008$ & $0.035\pm0.004$\\
\midrule
Hard-scatter finding efficiency & $\pm0.1$ mm & $99.90^{+0.07}_{-0.14}\%$ & $99.79^{+0.10}_{-0.16}\%$\\
\bottomrule
\end{tabular}
\par\vspace{.25em}\raggedright\footnotesize
\textbf{Efficiency:} truth PVs with $\ntrk\geq2$ selected reco tracks and a reco vertex within $0.1$ mm / all such truth PVs. One reco vertex can be near several truth PVs.\\[.3em]
\textbf{Unmatched:} no eligible truth PV within $0.1$ mm. \textbf{Extra nearby reco:} another candidate near an already matched truth PV.\\[.3em]
\textbf{Hard scatter:} events with any vertex within $0.1$ mm of the true HS position / all events.
\sourcecaption{HS follows the ITk paper's written $z$ criterion (Sec.\ 5.5, p.\ 27); the inclusive extension is our diagnostic. PV-Finder: peaks before GNN; AMVF: $\ntrk\geq2$. Unmatched reco may be near one-track truth PVs. Errors: 10,000 event bootstraps; HS: 68.27\% Clopper--Pearson.}
\end{frame}

\begin{frame}{One set of settings}
\vspace{.8em}
\centering
{\Large Integral threshold: \textbf{\textcolor{pennred}{0.3}}}\par
\vspace{1.8em}
{\Large Vertex finding: \textbf{\textcolor{pennred}{$|\Delta z|\leq0.1$ mm}}}\par
\vspace{1.8em}
{\Large Vertex categories: \textbf{\textcolor{pennred}{ATLAS-inspired diagnostic}}}\par
\vspace{.45em}
{\small Clean / Matched, Merged, Split, Fake}\par
\vfill
\sourcecaption{0.5 is a cross-check only. The all-vertex position metric is our diagnostic; the ITk paper uses 0.1 mm for HS finding. Current categories use unit track weights and an unlabelled-track proxy for Fake.}
\end{frame}

\begin{frame}{HL-LHC: vertex separation at integral cuts 0.3 and 0.5}
\centering
\includegraphics[width=.99\linewidth,height=.78\textheight,keepaspectratio]{figures/resolution_integral_030_050.pdf}
\sourcecaption{AMVF note Sec.\ 7.1, p.\ 14 shows the same kind of separation distribution; the fitted scale is our diagnostic. No truth-matching window enters this fit.}
\end{frame}

% Four independent examples of the agreed unweighted ntuple diagnostic.
% Every dot is one assigned track. The Split example includes both vertices;
% Fake is the unlabelled-track proxy, shown as a strict plurality below 50%.
\begin{frame}{Unweighted vertex categories: four examples}
\context{Each dot is one assigned track. Four independent examples of the convention used in our plots.}
\centering
\resizebox{.98\linewidth}{!}{%
\begin{tikzpicture}[x=1cm,y=1cm,font=\small]
  \colorlet{catA}{mydarkblue!85}
  \colorlet{catB}{myorange!85!black}
  \colorlet{catC}{teal!70!black}
  \colorlet{catF}{black!45}
  \foreach \x/\w in {1.6/3.0,5.0/3.0,9.35/5.0,13.7/3.0}{
    \node[pvcard,minimum width=\w cm,minimum height=3.65cm] at (\x,2.05) {};
  }
  % Clean: one truth interaction contributes 5/6 of the assigned tracks.
  \node[font=\small\bfseries,text=mygreen!65!black] at (1.6,3.5) {Clean / Matched};
  \foreach \p in {(1.0,2.75),(1.6,2.75),(2.2,2.75),(1.0,2.12),(1.6,2.12)}
    \fill[catA] \p circle (5pt);
  \fill[catB] (2.2,2.12) circle (5pt);
  \node[align=center,font=\footnotesize] at (1.6,1.1) {$5/6=83\%$ from one\\truth collision};
  % Merged: 60/40 avoids an ambiguous tie for the dominant truth origin.
  \node[font=\small\bfseries,text=myorange!85!black] at (5.0,3.5) {Merged};
  \foreach \p in {(4.4,2.75),(5.0,2.75),(5.6,2.75)}
    \fill[catA] \p circle (5pt);
  \foreach \p in {(4.7,2.12),(5.3,2.12)}
    \fill[catB] \p circle (5pt);
  \node[align=center,font=\footnotesize] at (5.0,1.1) {$60\%+40\%$\\Neither reaches $70\%$};
  % Split: both vertices are dominated by the same truth interaction.
  \node[font=\small\bfseries,text=myblue] at (9.35,3.5) {Split};
  \foreach \x/\name in {8.15/v_a,10.55/v_b}{
    \node[draw=black!25,rounded corners=3pt,minimum width=1.7cm,minimum height=1.1cm] at (\x,2.55) {};
    \node[font=\scriptsize,text=black!60] at (\x,3.08) {$\name$};
    \foreach \dx in {-.42,0,.42} \fill[catC] (\x+\dx,2.55) circle (5pt);
  }
  \node[font=\footnotesize] at (9.35,1.65) {$\sum p_T^2(v_a)>\sum p_T^2(v_b)$};
  \node[font=\footnotesize,text=mygreen!65!black] at (8.15,.98) {Keep Clean};
  \node[font=\footnotesize\bfseries,text=myblue] at (10.55,.98) {Relabel Split};
  \draw[-{Latex[length=2mm]},draw=myblue!65] (10.55,1.95)--(10.55,1.24);
  % Fake proxy: 3 unlabelled + 2 from each truth collision.
  \node[font=\small\bfseries,text=myred] at (13.7,3.5) {Fake (proxy)};
  \foreach \p in {(13.1,2.75),(13.7,2.75),(14.3,2.75)}
    \fill[catF] \p circle (5pt);
  \foreach \p in {(13.05,2.12),(13.48,2.12)} \fill[catA] \p circle (5pt);
  \foreach \p in {(13.92,2.12),(14.35,2.12)} \fill[catB] \p circle (5pt);
  \node[align=center,font=\footnotesize] at (13.7,1.1) {$3/7$ unlabelled $>2/7$\\from either collision};
  % Colours refer to truth origin within each independent example.
  \fill[catA] (1.0,-.45) circle (5pt);
  \node[anchor=west,font=\footnotesize] at (1.25,-.45) {truth collision A};
  \fill[catB] (4.7,-.45) circle (5pt);
  \node[anchor=west,font=\footnotesize] at (4.95,-.45) {truth collision B};
  \fill[catC] (8.4,-.45) circle (5pt);
  \node[anchor=west,font=\footnotesize] at (8.65,-.45) {same truth collision};
  \fill[catF] (12.65,-.45) circle (5pt);
  \node[anchor=west,font=\footnotesize] at (12.9,-.45) {unlabelled tracks};
\end{tikzpicture}}
\par\vspace{.65em}\small
Track counts set the category; $\sum p_T^2$ only resolves the Split label.
\sourcecaption{Unweighted approximation to \amvfnote. The paper uses fit weights and truth-classified fake tracks; neither is available here. Main results: top-1 GNN assignment, no score cut, reco $\ntrk\geq2$.}
\end{frame}


\begin{frame}{Vertex categories: our unweighted definitions}
\context{Every assigned track counts once. Truth labels identify its collision of origin.}
\small\renewcommand{\arraystretch}{1.5}
\begin{tabular}{p{3.2cm}p{10.3cm}}
\toprule
\textbf{Clean / Matched} & At least 70\% of the assigned tracks come from one true collision.\\
\textbf{Merged} & A true collision is the largest contribution, but none reaches 70\% of the assigned tracks.\\
\textbf{Split} & Several reco vertices have the same dominant true collision. Only the vertex with the largest $\sum p_T^2$ keeps Clean or Merged.\\
\textbf{Fake (proxy)} & The unlabelled-track group is the largest contribution.\\
\bottomrule
\end{tabular}
\par\vspace{.6em}\footnotesize
Split relabels duplicates after classifying their track content. Both methods retain reco vertices with at least two tracks. All category denominators include unlabelled tracks.
\sourcecaption{Adapted from \amvfnote; Track and Vertex Reconstruction with the ATLAS Inner Detector, Sec.\ 6.5, pp.\ 44--45 (local draft). Those papers use fit weights and truth-classified fake tracks. Our ntuples support this unweighted approximation.}
\end{frame}


\begin{frame}{PV-Finder + GNN: track-association diagnostics}
\context{Same 1,920 events, $185\leq\mu\leq200$; integral cut $0.3$; no GNN score cut; reco $\ntrk\geq2$ for both.}
\centering\small\renewcommand{\arraystretch}{1.2}
\begin{tabular}{lcc}
\toprule
Measurement & PV-Finder + GNN & AMVF (stored links)\\
\midrule
Track efficiency & \gnnEfficiency & \amvfEfficiency\\
Assignment purity & \gnnPurity & \amvfPurity\\
Assigned fraction & \gnnAssigned & \amvfAssigned\\
\midrule
Clean / reconstructed vertices & \gnnCleanReco & \amvfCleanReco\\
\bottomrule
\end{tabular}
\par\vspace{.25em}\raggedright\footnotesize
\textbf{Truth-labelled tracks:} efficiency = correct / all input tracks; purity = correct / assigned links; assigned fraction = assigned / all input tracks.\\[.3em]
Correct: the target's dominant unweighted truth origin equals the track's origin. All truth origins enter, including one-track origins; efficiency counts each track once.\\[.3em]
Clean uses the unit-weight approximation, including Split relabelling.
\takeaway{Our diagnostic: \gnnGainEfficiency\ efficiency points and \gnnGainPurity\ purity points over stored AMVF links.}
\sourcecaption{Top-1 GNN target, no score cut; remove reco vertices with fewer than two tracks, with no reassignment. Paired event-bootstrap errors. This does not evaluate an ATLAS TTVA working point; AMVF note Sec.\ 7.5, p.\ 18 instead plots compatible HS and pile-up track yields.}
\end{frame}

\begin{frame}{Proposed paper structure}
\begin{columns}[T]
\begin{column}{.45\linewidth}
\small
\textbf{1\quad Introduction}\\[.45em]
Why vertex reconstruction matters at high pile-up.\\[1em]
\textbf{2\quad ATLAS detector}\\[.6em]
\textbf{3\quad Run 3 Inner Detector}\\[.6em]
\textbf{4\quad The ITk upgrade}\\[1em]
\textbf{5\quad Data and MC samples}\\[.45em]
Run 3 data and simulation; HL-LHC simulation.\\[.35em]
Track selection, pile-up and beam spot.
\end{column}
\begin{column}{.51\linewidth}
\small
\textbf{6\quad PV-Finder and the track-to-vertex GNN}\\[.45em]
Inputs, both architectures, training and chosen settings.\\[1em]
\textbf{7\quad Performance}\\[.45em]
7.1\quad Run 3 simulation\\[.3em]
7.2\quad Run 3 data\\[.3em]
7.3\quad HL-LHC simulation\\[.3em]
7.4\quad Timing studies\\[1em]
\textbf{8\quad Conclusion}\\[.45em]
Main improvements and remaining limitations.
\end{column}
\end{columns}
\takeaway{Describe both algorithms in Section 6; compare their performance with AMVF in Section 7.}
\sourcecaption{Section order follows the professor's outline, dated 23 September 2026.}
\end{frame}

\begin{frame}{What we will include in the paper}
\small
\textbf{Use an integral cut of 0.3 for the main results.}\\[.4em]
Use the same setting for vertex yields, resolution and the GNN input vertices.
\par\vspace{.6em}
\centering\footnotesize\renewcommand{\arraystretch}{1.2}
\begin{tabular}{p{2.1cm}p{4.35cm}p{6.35cm}}
\toprule
Study & Available & Add or update\\
\midrule
Run 3 MC & V4 has Run 2 MC validation. & Add representative Run 3 MC; finder and GNN against AMVF.\\
Run 3 data & Agreement with AMVF positions and track assignments. & Compare the same observables in data and MC, with matched pile-up.\\
HL-LHC MC & Positional efficiency; unweighted track metrics and vertex categories. & Add HS selection and HS position-resolution curves; compare the chosen metrics consistently.\\
Timing & Athena CPU and GPU measurements. & Compare CPU with CPU; show finder, GNN and full-chain times versus pile-up.\\
\bottomrule
\end{tabular}
\par\vspace{.5em}\raggedright\small
Use the same unweighted category definitions for both methods; state the difference from ATLAS.
\sourcecaption{Main GNN results use no score cut; the score scan is a separate diagnostic. Re-evaluate Run 2/3 at integral cut 0.3 before quoting updated performance.}
\end{frame}

\begin{frame}{Comparison with the ITk paper}
\centering\small\renewcommand{\arraystretch}{1.25}
\begin{tabular}{p{3.05cm}p{3.05cm}p{6.6cm}}
\toprule
Measurement & Published result & Comparison with our study\\
\midrule
HS finding, $0.1$ mm & About $99.5\%$ & PV-Finder $99.90\%$, AMVF $99.79\%$. Same $z$ cut; exact monitoring selection to check.\\
HS finding + selection & About $92$--$93\%$ & Stored AMVF primary: $93.0\%$. The old $80.7\%$ re-ranked the ntuple track lists.\\
HS position precision & ITk: about $10\,\mu$m; Run 3: about $20\,\mu$m & Add HS resolution. Our inclusive width and vertex-separation scale are different quantities.\\
\bottomrule
\end{tabular}
\par\vspace{.6em}\raggedright\small
The stored primary contains the largest number of true HS tracks in 1,785/1,920 events ($93.0^{+0.6}_{-0.6}\%$). Re-ranking changes the chosen vertex in 428 events.
\takeaway{The integrated AMVF result is numerically close; matching samples and selections remains necessary.}
\sourcecaption{\itkpaper, Figs.\ 26--28. Our ntuple uses RecoVertex\_type = 1. Confirm geometry and the exact monitoring configuration before claiming a reproduction.}
\end{frame}

\begin{frame}{Run 3 timing: the published comparison}
\begin{columns}[c]
\begin{column}{.43\linewidth}
\small
\textbf{ACTS AMVF: about 200 ms/event at $\langle\mu\rangle\simeq90$.}\\[.7em]
That point is from a 2018 machine-development run, outside the standard data-quality selection.\\[.7em]
Benchmark: AMD EPYC 7302, one thread per reconstruction job.
\end{column}
\begin{column}{.54\linewidth}
\centering
\includegraphics[page=22,trim={172bp 562bp 157bp 104bp},clip,width=.96\linewidth,height=.48\textheight,keepaspectratio]{research/run3_tracking.pdf}
\end{column}
\end{columns}
\par\vspace{.9em}\small
V4's approximately $9\times$ gain compares \textbf{PV-Finder + GNN on GPU} with \textbf{AMVF on CPU}, on different hardware from the paper (Xeon 4216 / Tesla T4).
\par\vspace{.55em}
For the paper, show a common benchmark and separate finder, GNN and full-chain times.
\sourcecaption{\runthreepaper\ and benchmark description in Sec.\ 4; V4 timing source: \treysource.}
\end{frame}
